\documentclass[11pt,a4paper]{article}
\usepackage[margin=2.2cm]{geometry}
\usepackage[italian]{babel}
\usepackage[utf8]{inputenc}
\usepackage[T1]{fontenc}
\usepackage{graphicx}
\usepackage{subcaption}
\usepackage{booktabs}
\usepackage{array}
\usepackage{xcolor}
\usepackage{hyperref}
\usepackage{float}
\usepackage{csvsimple}

\hypersetup{colorlinks=true,linkcolor=black,urlcolor=blue,citecolor=black}
\setlength{\parindent}{0pt}
\setlength{\parskip}{0.55em}
\renewcommand{\arraystretch}{1.22}
\definecolor{mvpls}{HTML}{636363}
\definecolor{breg}{HTML}{FA9FB5}
\definecolor{bsmooth}{HTML}{DD3497}
\definecolor{pbpls}{HTML}{49006A}
\definecolor{fpcr}{HTML}{31A354}
\definecolor{rOne}{HTML}{3182BD}

\newcommand{\reviewcomment}[2]{%
  \begin{quote}\small
  \textbf{#1}\par
  \emph{#2}
  \end{quote}
}
\newcommand{\reviewhighlight}[1]{\textcolor{blue!70!black}{\textbf{#1}}}
\newcommand{\unmetreviewhighlight}[1]{\textcolor{red!70!black}{\textbf{#1}}}
\newcommand{\modelswatch}[1]{\colorbox{#1}{\rule{1.1em}{0.8ex}}}
\newcommand{\modellegend}{%
  \begin{center}\footnotesize
  \begin{tabular}{@{}ll@{\qquad}ll@{\qquad}ll@{}}
  \modelswatch{mvpls} & MV-PLS &
  \modelswatch{breg} & B-fPLS (Regression Splines) &
  \modelswatch{bsmooth} & B-fPLS (Smoothing Splines) \\
  \modelswatch{pbpls} & PB-fPLS &
  \modelswatch{fpcr} & fPCR &
  \modelswatch{rOne} & R1-fPLS
  \end{tabular}
  \end{center}
}
\newcommand{\modellegendtwod}{%
  \begin{center}\footnotesize
  \begin{tabular}{@{}ll@{\qquad}ll@{\qquad}ll@{}}
  \modelswatch{mvpls} & MV-PLS &
  \modelswatch{fpcr} & fPCR &
  \modelswatch{rOne} & R1-fPLS
  \end{tabular}
  \end{center}
}
\newcommand{\gcvcomponentfigure}[2]{%
  \begin{figure}[H]
  \centering
  \begin{subfigure}[t]{0.48\linewidth}
  \centering\includegraphics[page=1,width=\linewidth]{#1}
  \caption{$h=1$.}
  \end{subfigure}\hfill
  \begin{subfigure}[t]{0.48\linewidth}
  \centering\includegraphics[page=2,width=\linewidth]{#1}
  \caption{$h=2$.}
  \end{subfigure}
  \par\vspace{0.35em}
  \begin{subfigure}[t]{0.48\linewidth}
  \centering\includegraphics[page=3,width=\linewidth]{#1}
  \caption{$h=3$.}
  \end{subfigure}\hfill
  \begin{subfigure}[t]{0.48\linewidth}
  \centering\includegraphics[page=4,width=\linewidth]{#1}
  \caption{$h=4$.}
  \end{subfigure}
  \caption{#2}
  \end{figure}
}
\newcommand{\figureplaceholder}[2]{%
  \begin{figure}[H]
  \centering
  \fbox{\parbox[c][2.7cm][c]{0.92\linewidth}{\centering
    \textit{Placeholder figura}\\[0.35em]
    #1}}
  \caption{#2}
  \end{figure}
}
\newcommand{\submittedpaperfigure}[2][]{%
  \IfFileExists{#2}{%
    \includegraphics[#1]{#2}%
  }{%
    \fbox{\parbox[c][3.4cm][c]{0.82\linewidth}{\centering
      \textbf{\textcolor{red!70!black}{Asset esterno mancante}}\\[0.4em]
      La figura estratta dal manoscritto non è inclusa nel repository.\\
      Nessun risultato sostitutivo è mostrato.\\[0.35em]
      \texttt{\detokenize{#2}}}}%
  }%
}
\newcommand{\gcvtable}[3]{%
  \begin{table}[H]
  \centering
  \scriptsize
  \resizebox{\linewidth}{!}{%
  \begin{tabular}{@{}c c c c c c c@{}}
  \toprule
  $h$ & $\lambda_{\mathrm{fisso}}$ & $\lambda_{\mathrm{GCV}}^{\mathrm{dir}}$ &
  $\lambda_{\mathrm{GCV}}^{\mathrm{load}}$ & #2 fisso & #2 GCV & #2 MV-PLS \\
  \midrule
  \csvreader[head to column names,late after line=\\]{#1}{}{%
    \csvcoli & $\csvcolii$ & $\csvcoliii$ & $\csvcoliv$ &
    \csvcolv & \csvcolvi & \csvcolvii
  }
  \bottomrule
  \end{tabular}
  }
  \caption{#3}
  \end{table}
}
\newcommand{\reviewdesigntable}[1]{%
  \begin{table}[H]
  \centering
  \small
  \begin{tabular}{p{0.18\linewidth}p{0.72\linewidth}}
  \toprule
  \textbf{test suite / test} & \textbf{Setting nel paper} \\
  \midrule
  \csvreader[separator=pipe,head to column names,late after line=\\]{#1}{}{%
    \texttt{\csvcoli} & \csvcolii
  }
  \bottomrule
  \end{tabular}
  \caption{Test della simulazione sottomessa e configurazioni attuali.}
  \end{table}
}
\newcommand{\reviewiterationtable}[1]{%
  {\footnotesize
  \begin{center}
  \begin{tabular}{lccccc}
  \toprule
  Metodo & $h=1$ & $h=2$ & $h=3$ & $h=4$ & Fit monotono \\
  \midrule
  \csvreader[head to column names,late after line=\\]{#1}{}{%
    \csvcoli & \csvcolii & \csvcoliii & \csvcoliv & \csvcolv & \csvcolvi
  }
  \bottomrule
  \end{tabular}
  \end{center}
  \par}
}
\newcommand{\statusbox}[3]{%
  \vspace{0.35em}
  \noindent\fcolorbox{#1}{#2}{\parbox{0.94\linewidth}{#3}}%
  \vspace{0.4em}
}
\newcommand{\goodbox}[1]{\statusbox{green!50!black}{green!10}{\textbf{Confermato.} #1}}
\newcommand{\notebox}[1]{\statusbox{black!55}{gray!10}{#1}}
\newcommand{\warningbox}[1]{\statusbox{red!55!black}{red!10}{\textbf{Problema.} #1}}
\newcommand{\decisionbox}[1]{\statusbox{black!55}{gray!10}{#1}}
\newcommand{\bluebox}[1]{\statusbox{blue!60!black}{blue!8}{#1}}
\title{fPLS simulation study\\Confronto con la versione sottomessa e test per la review}
\author{Pietro Donelli}
\date{12 agosto 2026}

\begin{document}
\maketitle

\section{Confronto con la versione sottomessa}

L'implementazione di fPLS è stata portata nell'ultima versione della libreria, sul
\href{https://github.com/fdaPDE/fdaPDE-cpp/tree/develop_fPLS}{branch \texttt{develop\_fPLS}
di fdaPDE-cpp}, a partire da \texttt{stable}. I test seguenti usano questa versione e sono
implementati nel
\href{https://github.com/donellipietro/fPLS-simulation-study/tree/stable}{repository
\texttt{fPLS-simulation-study}, branch \texttt{stable}}.

Le quattro simulazioni già nel paper mantengono lo stesso setting. Le modifiche comuni sono:
\begin{itemize}
  \item 30 ripetizioni per setting, gli stessi livelli di rumore del paper
  ($\mathrm{NSR}_X=3$, $\mathrm{NSR}_Y=0.1$) e quattro componenti;
  \item tolleranza $10^{-8}$ e \texttt{max\_iter = 1000} per tutti i metodi iterativi;
  \item GCV nativo riproducibile per R1-fPLS e fPCR (seed fisso e 100 probe vectors);
  \item griglia delle penalizzazioni $10^{-9},10^{-8},\ldots,10^2$ per R1-fPLS e
  PB-fPLS; per B-fPLS con smoothing splines resta la griglia densa originale
  $10^{-9},10^{-8.9},\ldots,10^2$;
  \item aggiunta di fPCR come benchmark;
  \item in 2D la procedura di generazione delle location è equivalente, ma non riproduce
  le coordinate estratte dall'implementazione R originaria. Per questo i
  risultati sono compatibili, non identici punto per punto.
\end{itemize}

\reviewdesigntable{review_tables/submitted_design.csv}

\subsection{\texttt{1D/vary\_k}: numero di nodi}

\subsubsection{Confronto quantitativo}

\begin{figure}[H]
\centering
{\small\textbf{Versione sottomessa (Fig.~3).}\par}
\submittedpaperfigure[width=0.78\linewidth]{paper_figure_crops/1d_vary_k.png}\par
\vspace{0.35em}
{\small\textbf{Output aggiornato.}\par}
\begin{subfigure}[t]{0.32\linewidth}
\centering\includegraphics[width=\linewidth]{../images/1D/vary_k/Y_reconstruction.pdf}
\caption{RMSE $Y$.}
\end{subfigure}\hfill
\begin{subfigure}[t]{0.32\linewidth}
\centering\includegraphics[width=\linewidth]{../images/1D/vary_k/X_reconstruction_locations.pdf}
\caption{RMSE $X$.}
\end{subfigure}\hfill
\begin{subfigure}[t]{0.32\linewidth}
\centering\includegraphics[width=\linewidth]{../images/1D/vary_k/Beta1_locations.pdf}
\caption{RMSE $\beta_1$.}
\end{subfigure}
\caption{Confronto completo sul test che varia il numero di nodi.}
\end{figure}

\goodbox{I risultati quantitativi sono equivalenti a quelli sottomessi.}

\notebox{fPCR (verde) è molto competitivo in questo setting. La
Sezione~\ref{sec:predictive-order} mostra uno scenario costruito per evidenziare la
differenza tra componenti ordinate per varianza e componenti predittive.}

\subsubsection{Confronto qualitativo}

\begin{figure}[H]
\centering
{\small\textbf{Versione sottomessa (Fig.~4).}\par}
\submittedpaperfigure[width=0.88\linewidth]{paper_figure_crops/1d_qualitative.png}\par
\vspace{0.45em}
{\small\textbf{Output aggiornato.}\par}
\includegraphics[width=0.88\linewidth]{../images/1D/vary_k/qualitative_beta.pdf}
\caption{Confronto qualitativo per \texttt{vary\_k}.}
\end{figure}

\notebox{Per la revisione terrei la figura qualitativa già presente nel paper, così non
serve rigenerare l'immagine.}

\subsection{\texttt{1D/vary\_n}: numero di soggetti}

\begin{figure}[H]
\centering
{\small\textbf{Versione sottomessa (Fig.~5).}\par}
\submittedpaperfigure[width=0.78\linewidth]{paper_figure_crops/1d_vary_n.png}\par
\vspace{0.35em}
{\small\textbf{Output aggiornato.}\par}
\begin{subfigure}[t]{0.32\linewidth}
\centering\includegraphics[width=\linewidth]{../images/1D/vary_n/Y_reconstruction.pdf}
\caption{RMSE $Y$.}
\end{subfigure}\hfill
\begin{subfigure}[t]{0.32\linewidth}
\centering\includegraphics[width=\linewidth]{../images/1D/vary_n/X_reconstruction_locations.pdf}
\caption{RMSE $X$.}
\end{subfigure}\hfill
\begin{subfigure}[t]{0.32\linewidth}
\centering\includegraphics[width=\linewidth]{../images/1D/vary_n/Beta1_locations.pdf}
\caption{RMSE $\beta_1$.}
\end{subfigure}
\caption{Confronto completo sul test che varia $n$.}
\end{figure}

\goodbox{I risultati quantitativi sono equivalenti a quelli della versione sottomessa.}

\subsection{\texttt{1D/vary\_m}: numero di location osservate}

\begin{figure}[H]
\centering
{\small\textbf{Versione sottomessa (Fig.~S2).}\par}
\submittedpaperfigure[width=0.78\linewidth]{paper_figure_crops/1d_vary_m.png}\par
\vspace{0.35em}
{\small\textbf{Output aggiornato.}\par}
\begin{subfigure}[t]{0.32\linewidth}
\centering\includegraphics[width=\linewidth]{../images/1D/vary_m/Y_reconstruction.pdf}
\caption{RMSE $Y$.}
\end{subfigure}\hfill
\begin{subfigure}[t]{0.32\linewidth}
\centering\includegraphics[width=\linewidth]{../images/1D/vary_m/X_reconstruction_locations.pdf}
\caption{RMSE $X$.}
\end{subfigure}\hfill
\begin{subfigure}[t]{0.32\linewidth}
\centering\includegraphics[width=\linewidth]{../images/1D/vary_m/Beta1_locations.pdf}
\caption{RMSE $\beta_1$.}
\end{subfigure}
\caption{Confronto completo sul test che varia $m$.}
\end{figure}

\goodbox{I risultati quantitativi sono equivalenti a quelli dei supplementary.}

\subsection{\texttt{2D/vary\_n}: numero di soggetti}

\begin{figure}[H]
\centering
{\small\textbf{Versione sottomessa (Fig.~6).}\par}
\submittedpaperfigure[width=0.82\linewidth]{paper_figure_crops/2d_qualitative.png}\par
\vspace{0.45em}
{\small\textbf{Output aggiornato.}\par}
\includegraphics[width=0.82\linewidth]{../images/2D/vary_n/beta_representative.pdf}
\caption{Confronto qualitativo in 2D. Le scale cromatiche sono fissate per riga usando il range della verità.}
\end{figure}

\goodbox{Le mappe qualitative sono equivalenti a quelle della versione sottomessa.}

\begin{figure}[H]
\centering
{\small\textbf{Versione sottomessa (Fig.~7).}\par}
\submittedpaperfigure[width=0.82\linewidth]{paper_figure_crops/2d_vary_n.png}\par
\vspace{0.35em}
{\small\textbf{Output aggiornato.}\par}
\begin{subfigure}[t]{0.32\linewidth}
\centering\includegraphics[width=\linewidth]{../images/2D/vary_n/X_reconstruction_locations.pdf}
\caption{RMSE $X$.}
\end{subfigure}\hfill
\begin{subfigure}[t]{0.32\linewidth}
\centering\includegraphics[width=\linewidth]{../images/2D/vary_n/Beta1_locations.pdf}
\caption{RMSE $\beta_1$.}
\end{subfigure}\hfill
\begin{subfigure}[t]{0.32\linewidth}
\centering\includegraphics[width=\linewidth]{../images/2D/vary_n/Beta2_locations.pdf}
\caption{RMSE $\beta_2$.}
\end{subfigure}
\caption{Confronto quantitativo completo in 2D.}
\end{figure}

\goodbox{I risultati quantitativi sono equivalenti a quelli della versione sottomessa.}

\section{Test aggiuntivi}

\subsection{Fit time nelle figure del paper}

\reviewcomment{Reviewer 3, comment 3}{
Because this is a submission to Statistics and Computing, I think the computational side should
be developed more explicitly. The paper explains the algorithmic structure and mentions
implementation in an open-source C++ library, which is a real strength. However, the manuscript
currently provides very little quantitative information on computational cost, scalability, or
numerical behavior. This is especially important in the neuroimaging application, where the
predictors live on a mesh with more than 36,000 nodes. At a minimum, I would like to see
\reviewhighlight{additional reporting on runtime and, if possible, memory usage for the main
experiments; number of iterations required for convergence; practical computational cost of
GCV-based tuning; sensitivity of computation to the basis dimension or mesh resolution; and
discussion of numerical stability when extracting several components.} Even a modest runtime
or scaling study would substantially improve the paper for this journal.}

\notebox{Come deciso, nelle figure quantitative già presenti nel paper la legend viene
spostata sotto i pannelli e lo spazio liberato in basso a destra viene usato per
il fit time. La composizione mostrata qui è ancora provvisoria: sistemo le figure
definitive una volta concordati i risultati da inserire. Le tabelle riportano la mediana delle
iterazioni per componente su tutti i livelli e le ripetizioni del test; ``100\%'' indica che tutte
le traiettorie osservate sono monotone. Il trattino indica una diagnostica non disponibile oppure
un fit non iterativo. Il fit time include preprocessing e tuning specifici, fit e post-processing
in memoria; esclude generazione dei dati, serializzazione e caricamento dei risultati, e
valutazione.}

\subsubsection{Fig.~3: \texttt{1D/vary\_k}}
\begin{figure}[H]
\centering
\begin{subfigure}[t]{0.48\linewidth}
\centering\includegraphics[width=\linewidth]{../images/1D/vary_k/Y_reconstruction.pdf}
\caption{RMSE $Y$.}
\end{subfigure}\hfill
\begin{subfigure}[t]{0.48\linewidth}
\centering\includegraphics[width=\linewidth]{../images/1D/vary_k/X_reconstruction_locations.pdf}
\caption{RMSE $X$.}
\end{subfigure}
\par\vspace{0.35em}
\begin{subfigure}[t]{0.48\linewidth}
\centering\includegraphics[width=\linewidth]{../images/1D/vary_k/Beta1_locations.pdf}
\caption{RMSE $\beta_1$.}
\end{subfigure}\hfill
\begin{subfigure}[t]{0.48\linewidth}
\centering\includegraphics[width=\linewidth]{../images/1D/vary_k/fit_seconds.pdf}
\caption{Fit time.}
\end{subfigure}
\par\vspace{0.15em}
\modellegend
\caption{Proposta per la Fig.~3: legend sotto i pannelli, fit time in basso a destra.}
\end{figure}

\reviewiterationtable{review_tables/1D_vary_k_iterations.csv}

\notebox{La figura riporta il fit time al variare di $k$. B-fPLS con smoothing splines è il
metodo più costoso: effettua il pre-smoothing separatamente per ciascuna osservazione e usa una
griglia di $\lambda$ più densa. R1-fPLS e fPCR sono più veloci di PB-fPLS e di B-fPLS con
smoothing splines, con un vantaggio che cresce all'aumentare di $k$; B-fPLS con regression
splines e MV-PLS restano più leggeri in una parte dell'intervallo.}

\subsubsection{Fig.~5: \texttt{1D/vary\_n}}
\begin{figure}[H]
\centering
\begin{subfigure}[t]{0.48\linewidth}
\centering\includegraphics[width=\linewidth]{../images/1D/vary_n/Y_reconstruction.pdf}
\caption{RMSE $Y$.}
\end{subfigure}\hfill
\begin{subfigure}[t]{0.48\linewidth}
\centering\includegraphics[width=\linewidth]{../images/1D/vary_n/X_reconstruction_locations.pdf}
\caption{RMSE $X$.}
\end{subfigure}
\par\vspace{0.35em}
\begin{subfigure}[t]{0.48\linewidth}
\centering\includegraphics[width=\linewidth]{../images/1D/vary_n/Beta1_locations.pdf}
\caption{RMSE $\beta_1$.}
\end{subfigure}\hfill
\begin{subfigure}[t]{0.48\linewidth}
\centering\includegraphics[width=\linewidth]{../images/1D/vary_n/fit_seconds.pdf}
\caption{Fit time.}
\end{subfigure}
\par\vspace{0.15em}
\modellegend
\caption{Proposta per la Fig.~5: legend sotto i pannelli, fit time in basso a destra.}
\end{figure}

\reviewiterationtable{review_tables/1D_vary_n_iterations.csv}

\notebox{Qui R1-fPLS e fPCR usano la discretizzazione con 51 nodi, mentre B-fPLS e
PB-fPLS usano 11 nodi. R1-fPLS rimane comunque più veloce di PB-fPLS e di B-fPLS con
smoothing splines, mentre B-fPLS con regression splines e MV-PLS restano più leggeri in
questo setting.}

\subsubsection{Fig.~S2: \texttt{1D/vary\_m}}
\begin{figure}[H]
\centering
\begin{subfigure}[t]{0.48\linewidth}
\centering\includegraphics[width=\linewidth]{../images/1D/vary_m/Y_reconstruction.pdf}
\caption{RMSE $Y$.}
\end{subfigure}\hfill
\begin{subfigure}[t]{0.48\linewidth}
\centering\includegraphics[width=\linewidth]{../images/1D/vary_m/X_reconstruction_locations.pdf}
\caption{RMSE $X$.}
\end{subfigure}
\par\vspace{0.35em}
\begin{subfigure}[t]{0.48\linewidth}
\centering\includegraphics[width=\linewidth]{../images/1D/vary_m/Beta1_locations.pdf}
\caption{RMSE $\beta_1$.}
\end{subfigure}\hfill
\begin{subfigure}[t]{0.48\linewidth}
\centering\includegraphics[width=\linewidth]{../images/1D/vary_m/fit_seconds.pdf}
\caption{Fit time.}
\end{subfigure}
\par\vspace{0.15em}
\modellegend
\caption{Proposta per la Fig.~S2: legend sotto i pannelli, fit time in basso a destra.}
\end{figure}

\reviewiterationtable{review_tables/1D_vary_m_iterations.csv}

\notebox{Anche in questo setting R1-fPLS e fPCR usano la discretizzazione con 51 nodi,
mentre B-fPLS e PB-fPLS usano 11 nodi. R1-fPLS è più veloce di PB-fPLS e di B-fPLS con
smoothing splines, mentre B-fPLS con regression splines e MV-PLS restano più leggeri.}

\subsubsection{Fig.~7: \texttt{2D/vary\_n}}
\begin{figure}[H]
\centering
\begin{subfigure}[t]{0.48\linewidth}
\centering\includegraphics[width=\linewidth]{../images/2D/vary_n/X_reconstruction_locations.pdf}
\caption{RMSE $X$.}
\end{subfigure}\hfill
\begin{subfigure}[t]{0.48\linewidth}
\centering\includegraphics[width=\linewidth]{../images/2D/vary_n/Beta1_locations.pdf}
\caption{RMSE $\beta_1$.}
\end{subfigure}
\par\vspace{0.35em}
\begin{subfigure}[t]{0.48\linewidth}
\centering\includegraphics[width=\linewidth]{../images/2D/vary_n/Beta2_locations.pdf}
\caption{RMSE $\beta_2$.}
\end{subfigure}\hfill
\begin{subfigure}[t]{0.48\linewidth}
\centering\includegraphics[width=\linewidth]{../images/2D/vary_n/fit_seconds.pdf}
\caption{Fit time.}
\end{subfigure}
\par\vspace{0.15em}
\modellegendtwod
\caption{Proposta per la Fig.~7: legend sotto i pannelli, fit time in basso a destra.}
\end{figure}

\reviewiterationtable{review_tables/2D_vary_n_iterations.csv}

\notebox{Nel confronto 2D il fit time di fPCR è inferiore a quello di R1-fPLS nel perimetro
misurato. fPCR stima i loading, mentre R1-fPLS stima anche le direzioni e seleziona via GCV un
$\lambda$ per entrambe le quantità a ogni componente.}


\subsection{\texttt{2D/vary\_k}: risoluzione della mesh}

\reviewcomment{Reviewer 3, comment 3}{
Because this is a submission to Statistics and Computing, I think the computational side should
be developed more explicitly. The paper explains the algorithmic structure and mentions
implementation in an open-source C++ library, which is a real strength. However, the manuscript
currently provides very little quantitative information on computational cost, scalability, or
numerical behavior. This is especially important in the neuroimaging application, where the
predictors live on a mesh with more than 36,000 nodes. At a minimum, I would like to see
\reviewhighlight{additional reporting on runtime and, if possible, memory usage for the main
experiments; number of iterations required for convergence; practical computational cost of
GCV-based tuning; sensitivity of computation to the basis dimension or mesh resolution; and
discussion of numerical stability when extracting several components.} Even a modest runtime
or scaling study would substantially improve the paper for this journal.}

\textbf{Test added.} Nel setting 2D sottomesso variamo il numero di nodi della mesh.
I valori sono $k\in\{25,100,400,900,1600,2500\}$; manteniamo $n=100$, $m=1600$ e gli stessi
livelli di rumore. Il test separa l'effetto della risoluzione della mesh da quello del numero di
punti osservati.

\begin{figure}[H]
\centering
\begin{subfigure}[t]{0.48\linewidth}
\centering\includegraphics[width=\linewidth]{../images/2D/vary_k/Y_reconstruction.pdf}
\caption{RMSE $Y$.}
\end{subfigure}\hfill
\begin{subfigure}[t]{0.48\linewidth}
\centering\includegraphics[width=\linewidth]{../images/2D/vary_k/X_reconstruction_locations.pdf}
\caption{RMSE $X$.}
\end{subfigure}
\par\vspace{0.35em}
\begin{subfigure}[t]{0.48\linewidth}
\centering\includegraphics[width=\linewidth]{../images/2D/vary_k/Beta1_locations.pdf}
\caption{RMSE $\beta_1$.}
\end{subfigure}\hfill
\begin{subfigure}[t]{0.48\linewidth}
\centering\includegraphics[width=\linewidth]{../images/2D/vary_k/Beta2_locations.pdf}
\caption{RMSE $\beta_2$.}
\end{subfigure}
\par\vspace{0.35em}
\begin{subfigure}[t]{0.48\linewidth}
\centering\includegraphics[width=\linewidth]{../images/2D/vary_k/fit_seconds.pdf}
\caption{Fit time.}
\end{subfigure}
\par\vspace{0.15em}
\modellegendtwod
\caption{\texttt{2D/vary\_k}: accuratezza e costo al variare dei nodi della mesh.}
\end{figure}

\reviewiterationtable{review_tables/2D_vary_k_iterations.csv}

\goodbox{Come atteso, le ricostruzioni non sono influenzate dalla risoluzione della mesh una
volta raggiunto un numero sufficiente di nodi per rappresentare i campi stimati. L'unico
andamento sistematico è quello dei fit time.}

\decisionbox{Quali risultati di questo test vogliamo riportare nel paper o nei supplementary?}

\subsection{\texttt{1D/vary\_NSRX} e \texttt{1D/vary\_NSRY}: robustezza al rumore}

\reviewcomment{Reviewer 3, comment 4}{
The simulation studies are generally well designed and do a good job showing that the proposed
method can outperform presmoothing-based alternatives in the settings considered. I also
appreciate that the supplement includes additional results for varying the number of observed
locations. Still, for a journal such as Statistics and Computing, I think the computational
validation should go further. First, the comparison set becomes rather limited in the
multidimensional setting, where the main benchmark is essentially multivariate PLS. I
understand that off-the-shelf competitors may not exist for exactly the same setting, but the
manuscript should either add more relevant baselines or explain more clearly why broader
comparisons are not feasible. \reviewhighlight{Second, I would encourage the authors to
include at least some additional settings that probe the method's robustness and computational
behavior, for example varying mesh resolution, increasing missingness or noise, varying the
number of extracted components, or examining sensitivity to the tuning strategy.} For this
journal, requests for additional simulation settings seem entirely reasonable.}

\textbf{Test added.} Scenari più facili e più difficili del setting del paper, variando
separatamente il rumore in $X$ e in $Y$.

\begin{figure}[H]
\centering
\begin{subfigure}[t]{0.32\linewidth}
\centering\includegraphics[width=\linewidth]{../images/1D/vary_NSRX/Y_reconstruction.pdf}
\caption{RMSE $Y$.}
\end{subfigure}\hfill
\begin{subfigure}[t]{0.32\linewidth}
\centering\includegraphics[width=\linewidth]{../images/1D/vary_NSRX/X_reconstruction_locations.pdf}
\caption{RMSE $X$.}
\end{subfigure}\hfill
\begin{subfigure}[t]{0.32\linewidth}
\centering\includegraphics[width=\linewidth]{../images/1D/vary_NSRX/Beta1_locations.pdf}
\caption{RMSE $\beta_1$.}
\end{subfigure}
\caption{\texttt{1D/vary\_NSRX}: tutte le metriche al variare del rumore in $X$.}
\end{figure}

\begin{figure}[H]
\centering
\begin{subfigure}[t]{0.32\linewidth}
\centering\includegraphics[width=\linewidth]{../images/1D/vary_NSRY/Y_reconstruction.pdf}
\caption{RMSE $Y$.}
\end{subfigure}\hfill
\begin{subfigure}[t]{0.32\linewidth}
\centering\includegraphics[width=\linewidth]{../images/1D/vary_NSRY/X_reconstruction_locations.pdf}
\caption{RMSE $X$.}
\end{subfigure}\hfill
\begin{subfigure}[t]{0.32\linewidth}
\centering\includegraphics[width=\linewidth]{../images/1D/vary_NSRY/Beta1_locations.pdf}
\caption{RMSE $\beta_1$.}
\end{subfigure}
\caption{\texttt{1D/vary\_NSRY}: tutte le metriche al variare del rumore in $Y$.}
\end{figure}

L'aumento di RMSE $X$ con $\mathrm{NSR}_X$ è atteso. Il fatto che $\beta_1$ dei metodi
regolarizzati resti stabile non impedisce a RMSE $Y$ di crescere: la predizione di $Y$ usa una
$X$ più rumorosa e quindi ne è influenzata. Il comportamento al variare di
$\mathrm{NSR}_Y$ è quello atteso.

\decisionbox{Quali risultati vogliamo aggiungere al paper o ai supplementaries?}

\subsection{\texttt{2D/vary\_NSRX} e \texttt{2D/vary\_NSRY}: robustezza al rumore}

\reviewcomment{Reviewer 3, comment 4}{
The simulation studies are generally well designed and do a good job showing that the proposed
method can outperform presmoothing-based alternatives in the settings considered. I also
appreciate that the supplement includes additional results for varying the number of observed
locations. Still, for a journal such as Statistics and Computing, I think the computational
validation should go further. First, the comparison set becomes rather limited in the
multidimensional setting, where the main benchmark is essentially multivariate PLS. I
understand that off-the-shelf competitors may not exist for exactly the same setting, but the
manuscript should either add more relevant baselines or explain more clearly why broader
comparisons are not feasible. \reviewhighlight{Second, I would encourage the authors to
include at least some additional settings that probe the method's robustness and computational
behavior, for example varying mesh resolution, increasing missingness or noise, varying the
number of extracted components, or examining sensitivity to the tuning strategy.} For this
journal, requests for additional simulation settings seem entirely reasonable.}

\textbf{Test added.} Stessa estensione in 2D, per verificare che la robustezza non sia un
effetto della sola simulazione 1D.

\begin{figure}[H]
\centering
\begin{subfigure}[t]{0.48\linewidth}
\centering\includegraphics[width=\linewidth]{../images/2D/vary_NSRX/Y_reconstruction.pdf}
\caption{RMSE $Y$.}
\end{subfigure}\hfill
\begin{subfigure}[t]{0.48\linewidth}
\centering\includegraphics[width=\linewidth]{../images/2D/vary_NSRX/X_reconstruction_locations.pdf}
\caption{RMSE $X$.}
\end{subfigure}
\par\vspace{0.35em}
\begin{subfigure}[t]{0.48\linewidth}
\centering\includegraphics[width=\linewidth]{../images/2D/vary_NSRX/Beta1_locations.pdf}
\caption{RMSE $\beta_1$.}
\end{subfigure}\hfill
\begin{subfigure}[t]{0.48\linewidth}
\centering\includegraphics[width=\linewidth]{../images/2D/vary_NSRX/Beta2_locations.pdf}
\caption{RMSE $\beta_2$.}
\end{subfigure}
\caption{\texttt{2D/vary\_NSRX}: tutte le metriche al variare del rumore in $X$.}
\end{figure}

\begin{figure}[H]
\centering
\begin{subfigure}[t]{0.48\linewidth}
\centering\includegraphics[width=\linewidth]{../images/2D/vary_NSRY/Y_reconstruction.pdf}
\caption{RMSE $Y$.}
\end{subfigure}\hfill
\begin{subfigure}[t]{0.48\linewidth}
\centering\includegraphics[width=\linewidth]{../images/2D/vary_NSRY/X_reconstruction_locations.pdf}
\caption{RMSE $X$.}
\end{subfigure}
\par\vspace{0.35em}
\begin{subfigure}[t]{0.48\linewidth}
\centering\includegraphics[width=\linewidth]{../images/2D/vary_NSRY/Beta1_locations.pdf}
\caption{RMSE $\beta_1$.}
\end{subfigure}\hfill
\begin{subfigure}[t]{0.48\linewidth}
\centering\includegraphics[width=\linewidth]{../images/2D/vary_NSRY/Beta2_locations.pdf}
\caption{RMSE $\beta_2$.}
\end{subfigure}
\caption{\texttt{2D/vary\_NSRY}: tutte le metriche al variare del rumore in $Y$.}
\end{figure}

\goodbox{Tutto il comportamento osservato in 1D si ritrova anche in 2D: stabilità della
stima regolarizzata al crescere del rumore in $X$, crescita degli errori di ricostruzione e
risposta attesa al crescere del rumore in $Y$.}

\decisionbox{Quali risultati vogliamo aggiungere al paper o ai supplementaries?}

\subsection{\texttt{GCV/sensitivity} e \texttt{GCV/sensitivity\_high\_NSRX}}

\reviewcomment{Reviewer 2, comment 6}{
Section S2.1 gives the GCV criterion for selecting smoothing parameters, but the
implementation details are not clear. \reviewhighlight{Please specify how GCV was carried
out in the simulations and the real-data application, including the search range or
optimization procedure for $\lambda$. A brief sensitivity analysis or additional simulation
assessing the performance of GCV-selected tuning parameters would also strengthen the paper.}}

\reviewcomment{Reviewer 3, comment 4}{
The simulation studies are generally well designed and do a good job showing that the proposed
method can outperform presmoothing-based alternatives in the settings considered. I also
appreciate that the supplement includes additional results for varying the number of observed
locations. Still, for a journal such as Statistics and Computing, I think the computational
validation should go further. First, the comparison set becomes rather limited in the
multidimensional setting, where the main benchmark is essentially multivariate PLS. I
understand that off-the-shelf competitors may not exist for exactly the same setting, but the
manuscript should either add more relevant baselines or explain more clearly why broader
comparisons are not feasible. \reviewhighlight{Second, I would encourage the authors to
include at least some additional settings that probe the method's robustness and computational
behavior, for example varying mesh resolution, increasing missingness or noise, varying the
number of extracted components, or examining sensitivity to the tuning strategy.} For this
journal, requests for additional simulation settings seem entirely reasonable.}

\textbf{Test added.} Selezione GCV confrontata con MV-PLS ($\lambda=0$) e con
penalizzazioni fissate e uguali per direzioni e loading. Il confronto è ripetuto nello scenario
standard ($\mathrm{NSR}_X=3$) e con rumore alto in $X$ ($\mathrm{NSR}_X=10$), mantenendo
$\mathrm{NSR}_Y=0.1$, $n=100$, $m=200$, $k=51$ e quattro componenti. La griglia è
$10^{-9},\ldots,10^{-2}$. GCV seleziona invece una penalizzazione distinta per direzioni e
loading a ogni componente, con seed 0 e 100 probe.

Ciascun pannello, indicizzato da $h$, confronta il modello troncato ai primi $h$ componenti con
la corrispondente verità simulata. In particolare, le ricostruzioni stimate
$\widehat X_h$ e $\widehat Y_h$ sono confrontate con le ricostruzioni noise-free di verità
$X_h^\star$ e $Y_h^\star$; la stima cumulativa $\widehat\beta_h$ è confrontata con
$\beta_h^\star$; infine, la direzione, il loading e lo score stimati del componente $h$ sono
confrontati con le rispettive quantità di verità. L'ordine sequenziale dei componenti è
mantenuto senza permutation matching.

Le metriche sono definite come segue. Per una matrice $A$ e la corrispondente verità $A^\star$,

\[
  \mathrm{RMSE}(A,A^\star) =
  \sqrt{\frac{1}{|A|}\|A-A^\star\|_F^2},
  \qquad
  \mathrm{cos}_F(A,A^\star) =
  \frac{\langle A,A^\star\rangle_F}{\|A\|_F\|A^\star\|_F},
\]

\[
  \mathrm{cos}_{\mathrm{abs}}(a,a^\star) =
  \frac{|a^\top a^\star|}{\|a\|_2\|a^\star\|_2}.
\]

Qui $|A|$ indica il numero di elementi di $A$.

Si usa la RMSE per $(\widehat X_h,X_h^\star)$ e $(\widehat Y_h,Y_h^\star)$, il coseno di
Frobenius con segno per $(\widehat\beta_h,\beta_h^\star)$ e il coseno in valore assoluto
per le direzioni, i loading e gli score di $X$ del componente $h$. L'output aggregato conserva
anche la shape-RMSE di $\beta_h$, allineata nel segno e normalizzata come nei test del paper;
nelle tabelle di sensibilità qui sotto si riporta il coseno con segno.

Nelle tabelle seguenti, per ogni componente il lambda fisso riportato è quello che produce la
mediana minima per le RMSE e la mediana massima per le cosine similarity. Le colonne dei lambda
GCV riportano tutti i valori selezionati nelle 30 repliche, con la frequenza empirica tra parentesi;
i valori delle metriche nelle ultime tre colonne sono mediane sulle 30 repliche. La colonna MV-PLS
corrisponde al fit con $\lambda=0$.

\subsubsection{Scenario standard}

\gcvcomponentfigure{../images/GCV/sensitivity/Y_reconstruction.pdf}{%
\texttt{GCV/sensitivity}: RMSE della ricostruzione di $Y_h$.}
\gcvtable{gcv_tables/sensitivity_Y_reconstruction.csv}{RMSE $Y_h$}{%
\texttt{GCV/sensitivity}: migliori lambda fissi e risultati mediani per la ricostruzione di $Y_h$.}

\gcvcomponentfigure{../images/GCV/sensitivity/X_reconstruction_locations.pdf}{%
\texttt{GCV/sensitivity}: RMSE della ricostruzione di $X_h$.}
\gcvtable{gcv_tables/sensitivity_X_reconstruction.csv}{RMSE $X_h$}{%
\texttt{GCV/sensitivity}: migliori lambda fissi e risultati mediani per la ricostruzione di $X_h$.}

\gcvcomponentfigure{../images/GCV/sensitivity/Beta_locations.pdf}{%
\texttt{GCV/sensitivity}: cosine similarity con segno di $\beta_h$.}
\gcvtable{gcv_tables/sensitivity_Beta.csv}{$\mathrm{cos}_F(\beta_h)$}{%
\texttt{GCV/sensitivity}: migliori lambda fissi e risultati mediani per la cosine similarity con segno di $\beta_h$.}

\gcvcomponentfigure{../images/GCV/sensitivity/X_direction_locations.pdf}{%
\texttt{GCV/sensitivity}: cosine similarity assoluta delle direzioni di $X$.}
\gcvtable{gcv_tables/sensitivity_X_direction.csv}{$\mathrm{cos}_{\mathrm{abs}}$}{%
\texttt{GCV/sensitivity}: migliori lambda fissi e risultati mediani per la cosine similarity assoluta delle direzioni di $X$.}

\gcvcomponentfigure{../images/GCV/sensitivity/X_loading_locations.pdf}{%
\texttt{GCV/sensitivity}: cosine similarity assoluta dei loading di $X$.}
\gcvtable{gcv_tables/sensitivity_X_loading.csv}{$\mathrm{cos}_{\mathrm{abs}}$}{%
\texttt{GCV/sensitivity}: migliori lambda fissi e risultati mediani per la cosine similarity assoluta dei loading di $X$.}

\gcvcomponentfigure{../images/GCV/sensitivity/X_latent_scores.pdf}{%
\texttt{GCV/sensitivity}: cosine similarity assoluta degli score di $X$.}
\gcvtable{gcv_tables/sensitivity_X_scores.csv}{$\mathrm{cos}_{\mathrm{abs}}$}{%
\texttt{GCV/sensitivity}: migliori lambda fissi e risultati mediani per la cosine similarity assoluta degli score di $X$.}

\begin{figure}[H]
\centering
\begin{subfigure}[t]{0.48\linewidth}
\centering\includegraphics[width=\linewidth]{../images/GCV/sensitivity/fit_seconds.pdf}
\caption{Fit time.}
\end{subfigure}\hfill
\begin{subfigure}[t]{0.48\linewidth}
\centering\includegraphics[width=\linewidth]{../images/GCV/sensitivity/peak_ram_mib.pdf}
\caption{Peak RAM.}
\end{subfigure}
\caption{\texttt{GCV/sensitivity}: costo dell'intero fit a quattro componenti.}
\end{figure}

\subsubsection{Rumore alto in $X$}

\gcvcomponentfigure{../images/GCV/sensitivity_high_NSRX/Y_reconstruction.pdf}{%
\texttt{GCV/sensitivity\_high\_NSRX}: RMSE della ricostruzione di $Y_h$.}
\gcvtable{gcv_tables/sensitivity_high_NSRX_Y_reconstruction.csv}{RMSE $Y_h$}{%
\texttt{GCV/sensitivity\_high\_NSRX}: migliori lambda fissi e risultati mediani per la ricostruzione di $Y_h$.}

\gcvcomponentfigure{../images/GCV/sensitivity_high_NSRX/X_reconstruction_locations.pdf}{%
\texttt{GCV/sensitivity\_high\_NSRX}: RMSE della ricostruzione di $X_h$.}
\gcvtable{gcv_tables/sensitivity_high_NSRX_X_reconstruction.csv}{RMSE $X_h$}{%
\texttt{GCV/sensitivity\_high\_NSRX}: migliori lambda fissi e risultati mediani per la ricostruzione di $X_h$.}

\gcvcomponentfigure{../images/GCV/sensitivity_high_NSRX/Beta_locations.pdf}{%
\texttt{GCV/sensitivity\_high\_NSRX}: cosine similarity con segno di $\beta_h$.}
\gcvtable{gcv_tables/sensitivity_high_NSRX_Beta.csv}{$\mathrm{cos}_F(\beta_h)$}{%
\texttt{GCV/sensitivity\_high\_NSRX}: migliori lambda fissi e risultati mediani per la cosine similarity con segno di $\beta_h$.}

\gcvcomponentfigure{../images/GCV/sensitivity_high_NSRX/X_direction_locations.pdf}{%
\texttt{GCV/sensitivity\_high\_NSRX}: cosine similarity assoluta delle direzioni di $X$.}
\gcvtable{gcv_tables/sensitivity_high_NSRX_X_direction.csv}{$\mathrm{cos}_{\mathrm{abs}}$}{%
\texttt{GCV/sensitivity\_high\_NSRX}: migliori lambda fissi e risultati mediani per la cosine similarity assoluta delle direzioni di $X$.}

\gcvcomponentfigure{../images/GCV/sensitivity_high_NSRX/X_loading_locations.pdf}{%
\texttt{GCV/sensitivity\_high\_NSRX}: cosine similarity assoluta dei loading di $X$.}
\gcvtable{gcv_tables/sensitivity_high_NSRX_X_loading.csv}{$\mathrm{cos}_{\mathrm{abs}}$}{%
\texttt{GCV/sensitivity\_high\_NSRX}: migliori lambda fissi e risultati mediani per la cosine similarity assoluta dei loading di $X$.}

\gcvcomponentfigure{../images/GCV/sensitivity_high_NSRX/X_latent_scores.pdf}{%
\texttt{GCV/sensitivity\_high\_NSRX}: cosine similarity assoluta degli score di $X$.}
\gcvtable{gcv_tables/sensitivity_high_NSRX_X_scores.csv}{$\mathrm{cos}_{\mathrm{abs}}$}{%
\texttt{GCV/sensitivity\_high\_NSRX}: migliori lambda fissi e risultati mediani per la cosine similarity assoluta degli score di $X$.}

\begin{figure}[H]
\centering
\begin{subfigure}[t]{0.48\linewidth}
\centering\includegraphics[width=\linewidth]{../images/GCV/sensitivity_high_NSRX/fit_seconds.pdf}
\caption{Fit time.}
\end{subfigure}\hfill
\begin{subfigure}[t]{0.48\linewidth}
\centering\includegraphics[width=\linewidth]{../images/GCV/sensitivity_high_NSRX/peak_ram_mib.pdf}
\caption{Peak RAM.}
\end{subfigure}
\caption{\texttt{GCV/sensitivity\_high\_NSRX}: costo dell'intero fit a quattro componenti.}
\end{figure}

\goodbox{GCV seleziona quasi sempre $10^{-4}$ per direzioni e loading nello scenario
standard e con rumore alto; per la quarta direzione con $\mathrm{NSR}_X=10$ la mediana si
sposta a $10^{-3}$. Nelle ricostruzioni e nelle cosine similarity GCV segue da vicino il
miglior valore fissato della griglia in tutti i prefissi. Il vantaggio rispetto a MV-PLS cresce
nei componenti successivi ed è particolarmente evidente per $X$, $\beta$, direzioni, loading e
score quando il rumore in $X$ è alto.}

\notebox{Sul MacBook usato per questi risultati, la mediana del fit time è circa
$0.0163$ s per GCV, $0.0110$ s per MV-PLS e $0.0034$ s per un singolo fit nativo a
$\lambda$ fissato, in entrambi gli scenari. GCV valuta 64 criteri (otto valori, due
penalizzazioni e quattro componenti), inclusi i 100 probe. Il confronto GCV--fixed è diretto
perché usa lo stesso driver e lo stesso perimetro temporale; il confronto assoluto con MV-PLS
va interpretato come tempo dell'implementazione sulla macchina usata, perché MV-PLS è eseguito
in R su dati già in memoria.}

\decisionbox{Riportare i risultati di \texttt{sensitivity} e
\texttt{sensitivity\_high\_NSRX}.}

\subsection{\texttt{GCV/lambda\_to\_zero} e test a $\lambda$ molto piccolo}

\reviewcomment{Reviewer 2, comment 6}{
Section S2.1 gives the GCV criterion for selecting smoothing parameters, but the
implementation details are not clear. \reviewhighlight{Please specify how GCV was carried
out in the simulations and the real-data application, including the search range or
optimization procedure for $\lambda$. A brief sensitivity analysis or additional simulation
assessing the performance of GCV-selected tuning parameters would also strengthen the paper.}}

\reviewcomment{Reviewer 3, comment 3}{
Because this is a submission to Statistics and Computing, I think the computational side should
be developed more explicitly. The paper explains the algorithmic structure and mentions
implementation in an open-source C++ library, which is a real strength. However, the manuscript
currently provides very little quantitative information on computational cost, scalability, or
numerical behavior. This is especially important in the neuroimaging application, where the
predictors live on a mesh with more than 36,000 nodes. At a minimum, I would like to see
additional reporting on runtime and, if possible, memory usage for the main experiments; number
of iterations required for convergence; practical computational cost of GCV-based tuning;
sensitivity of computation to the basis dimension or mesh resolution; and
\reviewhighlight{discussion of numerical stability when extracting several components.} Even
a modest runtime or scaling study would substantially improve the paper for this journal.}

\textbf{Test added.} Comportamento della selezione GCV al crescere di $n$ e stabilità
numerica quando $\lambda$ diventa estremamente piccolo, nei casi $m>k$ e $m<k$.

\IfFileExists{../images/GCV/lambda_to_zero/Y_reconstruction.pdf}{%
\begin{figure}[H]
\centering
\begin{subfigure}[t]{0.32\linewidth}
\centering\includegraphics[width=\linewidth]{../images/GCV/lambda_to_zero/Y_reconstruction.pdf}
\caption{RMSE $Y$.}
\end{subfigure}\hfill
\begin{subfigure}[t]{0.32\linewidth}
\centering\includegraphics[width=\linewidth]{../images/GCV/lambda_to_zero/X_reconstruction_locations.pdf}
\caption{RMSE $X$.}
\end{subfigure}\hfill
\begin{subfigure}[t]{0.32\linewidth}
\centering\includegraphics[width=\linewidth]{../images/GCV/lambda_to_zero/Beta1_locations.pdf}
\caption{RMSE $\beta_1$.}
\end{subfigure}
\caption{\texttt{GCV/lambda\_to\_zero}: performance al crescere di $n$.}
\end{figure}

\begin{figure}[H]
\centering
\begin{subfigure}[t]{0.48\linewidth}
\centering\includegraphics[width=\linewidth]{../images/GCV/lambda_to_zero/direction_lambda.pdf}
\caption{Lambda delle direzioni.}
\end{subfigure}\hfill
\begin{subfigure}[t]{0.48\linewidth}
\centering\includegraphics[width=\linewidth]{../images/GCV/lambda_to_zero/loading_lambda.pdf}
\caption{Lambda dei loading.}
\end{subfigure}
\caption{\texttt{GCV/lambda\_to\_zero}: penalizzazioni selezionate al crescere di $n$.}
\end{figure}

\warningbox{Il test non mostra una convergenza semplice e generale di GCV a zero nel range
osservato. Non supporta una claim sulla convergenza di GCV a zero.}

\begin{figure}[H]
\centering
\begin{subfigure}[t]{0.32\linewidth}
\centering\includegraphics[width=\linewidth]{../images/GCV/small_lambda_m_gt_k/finite_solution.pdf}
\caption{Soluzioni finite.}
\end{subfigure}\hfill
\begin{subfigure}[t]{0.32\linewidth}
\centering\includegraphics[width=\linewidth]{../images/GCV/small_lambda_m_gt_k/hit_max_iter.pdf}
\caption{Max iteration raggiunto.}
\end{subfigure}\hfill
\begin{subfigure}[t]{0.32\linewidth}
\centering\includegraphics[width=\linewidth]{../images/GCV/small_lambda_m_gt_k/fit_seconds.pdf}
\caption{Fit time.}
\end{subfigure}
\caption{Test a $\lambda$ molto piccolo quando $m>k$.}
\end{figure}

\begin{figure}[H]
\centering
\begin{subfigure}[t]{0.32\linewidth}
\centering\includegraphics[width=\linewidth]{../images/GCV/small_lambda_m_lt_k/finite_solution.pdf}
\caption{Soluzioni finite.}
\end{subfigure}\hfill
\begin{subfigure}[t]{0.32\linewidth}
\centering\includegraphics[width=\linewidth]{../images/GCV/small_lambda_m_lt_k/hit_max_iter.pdf}
\caption{Max iteration raggiunto.}
\end{subfigure}\hfill
\begin{subfigure}[t]{0.32\linewidth}
\centering\includegraphics[width=\linewidth]{../images/GCV/small_lambda_m_lt_k/fit_seconds.pdf}
\caption{Fit time.}
\end{subfigure}
\caption{Test a $\lambda$ molto piccolo quando $m<k$.}
\end{figure}
}{%
\figureplaceholder{\texttt{GCV/lambda\_to\_zero}: performance al crescere di $n$.}{%
Performance al crescere di $n$.}
\figureplaceholder{\texttt{GCV/lambda\_to\_zero}: penalizzazioni selezionate al crescere di $n$.}{%
Penalizzazioni selezionate al crescere di $n$.}
\figureplaceholder{Test a $\lambda$ molto piccolo quando $m>k$.}{%
Test a $\lambda$ molto piccolo quando $m>k$.}
\figureplaceholder{Test a $\lambda$ molto piccolo quando $m<k$.}{%
Test a $\lambda$ molto piccolo quando $m<k$.}
}

\goodbox{I risultati a $\lambda$ molto piccolo sono quelli attesi. Per $m>k$ tutti i fit restano
finiti fino a $10^{-300}$. Per $m<k$ il test è interamente finito da $10^{-42}$ in poi e la
failure è gestita per valori più piccoli, quando il problema diventa numericamente non finito.}

\decisionbox{Aggiungere nei supplementary: il solo test di stabilità a lambda piccoli,
includendo la failure gestita per valori non finiti.}

\subsection{\texttt{2D/vary\_H}: numero di componenti estratte}

\reviewcomment{Reviewer 2, comment 7}{
\reviewhighlight{After Algorithm 1, the authors may clarify how the number of latent
components $H$ is selected.} Similar to the choice of the number of principal components in
PCA, $H$ is an important tuning parameter in PLS-type methods and may affect both prediction
and interpretation. In the simulation studies, it is not clear whether $H$ is fixed to the
true number of latent components used in data generation, or selected by a data-driven
criterion. In the real-data application, the authors retain $H=3$, but the selection criterion
is again not fully specified. I suggest adding a brief discussion of possible criteria for
choosing $H$, such as cross-validation or explained covariance, and reporting how $H$ was
chosen in the simulations and application.}

\reviewcomment{Reviewer 3, comment 4}{
The simulation studies are generally well designed and do a good job showing that the proposed
method can outperform presmoothing-based alternatives in the settings considered. I also
appreciate that the supplement includes additional results for varying the number of observed
locations. Still, for a journal such as Statistics and Computing, I think the computational
validation should go further. First, the comparison set becomes rather limited in the
multidimensional setting, where the main benchmark is essentially multivariate PLS. I
understand that off-the-shelf competitors may not exist for exactly the same setting, but the
manuscript should either add more relevant baselines or explain more clearly why broader
comparisons are not feasible. \reviewhighlight{Second, I would encourage the authors to
include at least some additional settings that probe the method's robustness and computational
behavior, for example varying mesh resolution, increasing missingness or noise, varying the
number of extracted components, or examining sensitivity to the tuning strategy.} For this
journal, requests for additional simulation settings seem entirely reasonable.}

\textbf{Test added.} Un unico fit fino a $H=5$, valutato progressivamente per separare il
comportamento empirico dalla scelta di $H$ nel data analysis.

\begin{figure}[H]
\centering
\begin{subfigure}[t]{0.48\linewidth}
\centering\includegraphics[width=\linewidth]{../images/2D/vary_H/Y_reconstruction.pdf}
\caption{Ricostruzione di $Y$ per numero di componenti.}
\end{subfigure}\hfill
\begin{subfigure}[t]{0.48\linewidth}
\centering\includegraphics[width=\linewidth]{../images/2D/vary_H/X_reconstruction_locations.pdf}
\caption{Ricostruzione di $X$ per numero di componenti.}
\end{subfigure}
\par\vspace{0.35em}
\begin{subfigure}[t]{0.48\linewidth}
\centering\includegraphics[width=\linewidth]{../images/2D/vary_H/Beta1_locations.pdf}
\caption{RMSE $\beta_1$ per numero di componenti.}
\end{subfigure}\hfill
\begin{subfigure}[t]{0.48\linewidth}
\centering\includegraphics[width=\linewidth]{../images/2D/vary_H/Beta2_locations.pdf}
\caption{RMSE $\beta_2$ per numero di componenti.}
\end{subfigure}
\caption{Effetto del numero di componenti nel test 2D: tutte le metriche.}
\end{figure}

\notebox{Da questi risultati inizia a vedersi la differenza tra fPCR e fPLS nella stima delle
componenti. fPCR ordina le componenti massimizzando la varianza spiegata di $X$: la
ricostruzione di $X$ progredisce quindi più rapidamente, mentre quella di $Y$ rimane indietro.
Nella generazione, i loading delle quattro componenti nei due response sono
\[
  C = \left[\begin{array}{rrrr}
  3 & 1 & 0 & 0 \\
  -3 & 0 & 1 & 4
  \end{array}\right],
\]
con righe corrispondenti a $Y_1,Y_2$. La ricostruzione di $\beta_2$ dipende quindi anche dalla
quarta componente, il cui loading in $Y_2$ è $4$. Stimare una quinta componente, oltre alle
quattro di segnale, peggiora la ricostruzione delle $\beta$, mentre l'impatto sulla ricostruzione
di $X$ e $Y$ è minore.}

Le traiettorie seguenti sono calcolate sui dati osservati, senza usare la verità simulata. Per
ogni numero $h$ di componenti, le quote di varianza ricostruita sono
\[
R_X^2(h) = 1 - \frac{\|X - \widehat X_h\|_F^2}{\|X - \bar X\|_F^2},
\qquad
R_Y^2(h) = 1 - \frac{\|Y - \widehat Y_h\|_F^2}{\|Y - \bar Y\|_F^2},
\]
dove $\bar X$ e $\bar Y$ sono le medie colonnari dei dati osservati. Per la
cross-covarianza riportiamo invece
\[
R_{XY}^2(h) = 1 -
\frac{\|C_{XY} - \widehat C_{XY,h}\|_F^2}{\|C_{XY}\|_F^2}.
\]
Potremmo proporre queste metriche per la selezione di $H$.

\begin{figure}[H]
\centering
\begin{subfigure}[t]{0.48\linewidth}
\centering\includegraphics[width=\linewidth]{../images/2D/vary_H/X_variance_explained.pdf}
\caption{Varianza osservata di $X$ ricostruita.}
\end{subfigure}\hfill
\begin{subfigure}[t]{0.48\linewidth}
\centering\includegraphics[width=\linewidth]{../images/2D/vary_H/Y_variance_explained.pdf}
\caption{Varianza osservata di $Y$ ricostruita.}
\end{subfigure}
\par\vspace{0.35em}
\begin{subfigure}[t]{0.72\linewidth}
\centering\includegraphics[width=\linewidth]{../images/2D/vary_H/XY_cross_covariance_explained.pdf}
\caption{Struttura di cross-covarianza osservata $X$--$Y$ ricostruita.}
\end{subfigure}
\modellegendtwod
\caption{Diagnostiche osservate lungo la traiettoria delle componenti nel test \texttt{2D/vary\_H}.}
\end{figure}

\notebox{fPCR ricostruisce una quota maggiore della varianza di $X$ con due e tre componenti,
mentre R1-fPLS e MV-PLS ricostruiscono una quota maggiore di $Y$. La cross-covarianza è già
in larga parte ricostruita con una componente e cresce verso il $100\%$ con le successive; è
informativa come controllo. MV-PLS tende invece a overfittare: alla quinta componente aumenta
la varianza osservata di $Y$ ricostruita, mentre i due metodi regolarizzati rimangono
sostanzialmente stabili.}

\decisionbox{Quali risultati di questo test vogliamo riportare nel paper o nei supplementary?}

\subsection{\texttt{2D-predictive-order/vary\_H}: componente predittiva a bassa varianza}
\label{sec:predictive-order}

\reviewcomment{Reviewer 2, comment 8}{
Since the introduction discusses the distinction between PLS and principal component regression,
and also refers to comparisons between fPLS and fPCA-based regression methods, it would be
useful to \reviewhighlight{include an fPCA-based functional principal component regression
method as an additional benchmark in the simulation study, at least for one setting considered
in simulation.} This comparison would help readers better assess the practical advantage of
R1-fPLS over another widely used functional dimension-reduction approach.}

\bluebox{fPCR è molto competitivo e nei test 1D non stiamo ancora giustificando in modo
esplicito il vantaggio di fPLS. Il test seguente è stato aggiunto per verificare se un setting
con componenti predittive a bassa varianza rendesse questa differenza visibile.}

\textbf{Test added.} A differenza di \texttt{2D/vary\_H}, che mantiene il setting sottomesso
con quattro componenti spaziali e la matrice dei loading di risposta $C$ riportata sopra, questo
è un setting sperimentale costruito per separare varianza di $X$ e rilevanza predittiva. Riutilizziamo
esattamente la data generation del test 2D sottomesso, estendendo da quattro a dieci le componenti
spaziali trigonometriche. Gli score latenti $\xi_1,\ldots,\xi_{10}$ sono generati come nel test principale,
con varianze di popolazione decrescenti; le prime otto componenti, più variabili, non entrano nella
risposta; poniamo invece
\[
Y_1 = \frac{\xi_9}{\sigma_9}, \qquad
Y_2 = \frac{\xi_{10}}{\sigma_{10}},
\]
così che le ultime due componenti, pur essendo le meno variabili in $X$, contribuiscano con la
stessa varianza a $Y$. Per tenere la difficoltà comparabile al resto degli esperimenti 2D, usiamo
il setting sottomesso: $n=100$, $m=1600$, mesh $30\times30$, $\mathrm{NSR}_X=3$ sulla
componente finale e $\mathrm{NSR}_Y=0.1$. Il test non è quindi una replica del setting principale:
modifica volutamente solo l'ordinamento delle componenti e le due colonne della matrice di risposta
associate alle componenti 9 e 10.
Prima di aggiungere il rumore verifichiamo inoltre, come nel test principale, che la decomposizione
PLS multivariata con tutte le componenti ricostruisca esattamente sia $X$ sia $Y$.

\begin{figure}[H]
\centering
\begin{subfigure}[t]{0.48\linewidth}
\centering\includegraphics[width=\linewidth]{../images/2D-predictive-order/vary_H/Y_reconstruction.pdf}
\caption{Ricostruzione di $Y$ al crescere di $H$.}
\end{subfigure}\hfill
\begin{subfigure}[t]{0.48\linewidth}
\centering\includegraphics[width=\linewidth]{../images/2D-predictive-order/vary_H/X_reconstruction_locations.pdf}
\caption{RMSE $X$.}
\end{subfigure}
\par\vspace{0.35em}
\begin{subfigure}[t]{0.48\linewidth}
\centering\includegraphics[width=\linewidth]{../images/2D-predictive-order/vary_H/Beta1_locations.pdf}
\caption{RMSE $\beta_1$.}
\end{subfigure}\hfill
\begin{subfigure}[t]{0.48\linewidth}
\centering\includegraphics[width=\linewidth]{../images/2D-predictive-order/vary_H/Beta2_locations.pdf}
\caption{RMSE $\beta_2$.}
\end{subfigure}
\caption{Test predittivo sperimentale per distinguere l'ordinamento fPCR e fPLS: tutte le metriche.}
\end{figure}

Le tre diagnostiche seguenti usano le stesse definizioni di $R_X^2(H)$, $R_Y^2(H)$ e
$R_{XY}^2(H)$ riportate nel test \texttt{2D/vary\_H}.

\begin{figure}[H]
\centering
\begin{subfigure}[t]{0.48\linewidth}
\centering\includegraphics[width=\linewidth]{../images/2D-predictive-order/vary_H/X_variance_explained.pdf}
\caption{Varianza osservata di $X$ ricostruita.}
\end{subfigure}\hfill
\begin{subfigure}[t]{0.48\linewidth}
\centering\includegraphics[width=\linewidth]{../images/2D-predictive-order/vary_H/Y_variance_explained.pdf}
\caption{Varianza osservata di $Y$ ricostruita.}
\end{subfigure}
\par\vspace{0.35em}
\begin{subfigure}[t]{0.72\linewidth}
\centering\includegraphics[width=\linewidth]{../images/2D-predictive-order/vary_H/XY_cross_covariance_explained.pdf}
\caption{Struttura di cross-covarianza osservata $X$--$Y$ ricostruita.}
\end{subfigure}
\modellegendtwod
\caption{Diagnostiche osservate nel test di ordinamento predittivo.}
\end{figure}

\warningbox{Il test non va interpretato come una giustificazione di $H=2$ per PLS. Sebbene
$Y$ sia generata da $\xi_9$ e $\xi_{10}$, queste sono componenti generative, non componenti PLS.
Le componenti PLS dipendono dalla covarianza campionaria di $X$ e dalla deflazione sequenziale.
Nella prima ripetizione del controllo noiseless la ricostruzione di $Y$ ha ancora RMSE $0.748$ con $H=2$, scende a $0.048$
con $H=7$ e diventa esatta solo con $H=10$.}

\goodbox{Il confronto rimane però informativo: PLS recupera il segnale di $Y$ molto prima di
fPCR. Con quattro componenti R1-fPLS e MV-PLS hanno già ricostruito gran parte della varianza
osservata di $Y$, e con sei componenti sono vicini al $90\%$; fPCR resta invece su quote molto
basse perché privilegia le componenti più variabili di $X$. fPCR recupera la prima direzione
predittiva alla componente 9 e la seconda alla 10.}

\decisionbox{Decidere se riportare questo test come evidenza qualitativa che PLS recupera il
segnale predittivo prima di fPCR. Se vogliamo sostenere la claim più forte ``due componenti contro
dieci'', il setting va ridisegnato imponendo una decomposizione PLS noiseless esatta a $H=2$.}

\notebox{In tutte le 30 ripetizioni del test, R1-fPLS e fPCR producono soluzioni finite e
traiettorie monotone; nessun fit raggiunge il limite di iterazioni.}

\section{Richieste dei reviewer ancora aperte}

Questa sezione riporta le richieste che restano aperte nel perimetro delle simulazioni. Le parti
in rosso non sono soddisfatte dai test eseguiti; le richieste sull'applicazione non sono riportate
qui.

\reviewcomment{Reviewer 2, comment 7}{
After Algorithm 1, the authors may clarify how the number of latent components $H$ is selected.
Similar to the choice of the number of principal components in PCA, $H$ is an important tuning
parameter in PLS-type methods and may affect both prediction and interpretation. In the simulation
studies, it is not clear whether $H$ is fixed to the true number of latent components used in data
generation, \unmetreviewhighlight{or selected by a data-driven criterion}. In the real-data
application, the authors retain $H=3$, but the selection criterion is again not fully specified. I
suggest adding a brief discussion of possible criteria for choosing $H$, such as
\unmetreviewhighlight{cross-validation or explained covariance, and reporting how $H$ was chosen
in the simulations} and application.}

\reviewcomment{Reviewer 3, comment 2}{
The proposed methodology is coherent and well motivated, and I find the regularized rank-one
construction appealing. However, the theoretical support for the actual penalized functional
estimator remains limited relative to the breadth of the claims. The supplement provides derivations
for the multivariate PLS characterization and for the operator representation, but I did not find
formal statistical guarantees for the proposed R1-fPLS estimator itself, such as consistency, error
bounds, or a clearer discussion of how discretization and regularization affect estimation accuracy.
I am not necessarily asking for a full asymptotic theory if that is beyond the intended scope of the
paper, but the authors should either provide stronger theoretical support or moderate the claims
accordingly.
In addition, several practical methodological choices need clearer discussion:
\unmetreviewhighlight{how the number of extracted components is selected, how sensitive the method
is to initialization}, what stopping criterion is used in practice, and
\unmetreviewhighlight{how the tuning of the smoothing parameters interacts with the sequential
extraction of components}. These details are important both statistically and computationally.}

\end{document}
